\chapter{Herramientas matemáticas}

En este capítulo se introducirán las matemáticas y geometría necesaria para comprender el resto de capítulos. Se verá una breve introducción a la geometría analítica y a vectores.

\section{Vectores}

En primer lugar se desarrollan conceptos relativos a \emph{espacios vectoriales}, estudiando sus características y propiedades. Esto será muy útil ya que proporcionará una estructura matemática para describir magnitudes como velocidad, posición, aceleración, entre otras; que serán especialmente útiles en la mayor parte del libro.

A continuación se presentarán geométrica y analíticamente los vectores en espacios bi y tridimensional, donde se dan las definiciones y las operaciones entre ellos, estableciendo las propiedades básicas.

\subsection{Espacios Vectoriales}

Para poder definir un espacio vectorial primero debe conocer la definición de vector,
\begin{definition}[Vector]
  Un vector \(\vec{v}\) en el plano xy es un par ordenado de números reales \((v_1,v_2)\). Los números \(v_1\) y \(v_2\) reciben el nombre de componentes del vector \(\vec{v}\).
\end{definition}

